%===================================================================================================
% 4. International taxation: majority support for taxing the rich to finance LICs
%===================================================================================================

\section{International taxation: majority support for taxing the rich to finance LICs}


%===================================================================================================
% - plausible policies: Fabre 26, Cappelen et al 2025
%===================================================================================================

%LB: Le texte qui suit est une paraphrase Fabre 26. Mais l'ensemble des 'plausible policies' n'apparait pas. Donc on pour'ait à la place de  'The 2\% tax on billionaire wealth is the most supported policy totalling 81\% acceptance [...]  gather overall, 70\% of acceptance.' écrire seulement le paragraphe d'après qui détaille l'ensemble des politiques et leur support en %.
\citet{fabre_public_2026} find in surveyed countries that all respondents accept\footnote{There is a distinction between \textit{support} and \textit{acceptance}. The term \textit{support} refers to the absolute proportion of \textit{Somewhat agree} or \textit{Strongly agree} responses on Likert scales, whereas the term \textit{acceptance} refers to relative support, that is, the proportion of support among responses non-\textit{indifferent}.} in majority global redistributive policies tested. The 2\% tax on billionaire wealth is the most supported policy totalling 81\% acceptance. Other proposals as debt reliefs for vulnerable countries -- where developed countries contribute at 0.7\% of the GDP in foreign aid --, expansion of the UN Security Council, or even, the Bridgetown initiative (an sustainable increase of investments at low interests rates in low income countries), gather overall, 70\% of acceptance. 

%LB: paragraphe 6.1. Figures 10, S23
%LB: remplacement du texte par l'image que celui ci décrivait. 
\citet{fabre_public_2026} tests support for ten plausible policies currently debated in international organizations. Figure \ref{fig:plausible_policies} reports their acceptance\footnote{There is a distinction between \textit{support} and \textit{acceptance}. The term \textit{support} refers to the absolute proportion of \textit{Somewhat agree} or \textit{Strongly agree} responses on Likert scales, whereas the term \textit{acceptance} refers to relative support, that is, the proportion of support among responses non-\textit{indifferent}.} across the surveyed countries. In the pooled sample, every proposal receives majority acceptance, ranging from 53\% for the international levy on aviation emissions to 81\% for the 2\% minimum tax on billionaires' wealth.

\begin{figure}[htbp]
    \centering
    \includegraphics[width=\textwidth]{figures/fabre_2026_figure_10.pdf}
    \caption{Acceptance of plausible global redistribution policies. \textit{Source:} \citet[Figure 10]{fabre_public_2026}.}
    \label{fig:plausible_policies}
\end{figure}

%LB: La catégorie U.S. voters for trump was at which date ? first or second mandate ? AF: second but no need to specify
The ``plausible'' label comes from the international debates upon these policies\footnote{As documented in \citet[Appendix C.1]{fabre_public_2026}, where polices as been discussed during summit as G20, during COPs or at ONU and OCDE}. As mentioned above, almost every policy proposal gather a majority of acceptance. Nevertheless, Japanese respondents acceptance on aviation taxation is at 46\% and it could be explained by the salient cost displayed through the 30\% increase in flight prices. For all plausible policies, left-leaning Europeans garner the highest acceptance rates (from 71\% for aviation taxation to 94\% for the 2\% billionaire taxation) followed by the U.S. Harris voters respondents (respectively, 67\% to 93\%). On the opposite, U.S. respondents who voted for Trump and far right voters in Europe display a reduced support compare to the overall mean (e.g., 34\% of Trump voters support aviation tax while 56\% support 2\% billionaire tax, the corresponding numbers are 40\% and 74\% for far-right Europeans voters) \citep[Figure S71]{fabre_public_2026}.

% Pour 'plausible policies' c'est plus délicat dans la papier de Capelen 2025 pour trouver quelque chose de concret. La seule information que l'on peut retenir que (i) l'universalisme des répondants et leur opinions politiques sont liées (donc universalism -> vote -> programme incluant intl' policies) (ii) les répondants donnent sur leur dotation, en moyenne, plus à des étrangers d'autres pays que des étrangers dans leur pays.
To complete these arguments: acceptance of plausible policies, \citet{cappelen_universalism_2025} ask respondents is the national government provide foreign aid in order to reduce ``economic differences between the rich and the poor'' in selected countries rather ``national government should focus on helping the poor in [country], rather than the poor elsewhere in the world''. Respondents generally prioritize the domestic poor: after responses are reverse-coded so that higher values indicate greater priority for the global poor, the mean lies around 1.5 on a four-point scale \citep[Figure 7B]{cappelen_universalism_2025}.
In the first experiment authors ask respondents to split hypothetical \$1000 between a in-group member (family, friend, neighbor, co-religionist, or co-ethnic) or a stranger in their country. The second experiment have the same method but oppose a stranger from her country to a stranger elsewhere in the world, respondents are explicitly told that the two recipients have the same standard of living. On average, respondents allocate 38.1\% to domestic strangers and 37.5\% to foreign strangers while 27\% divide the money equally in all decisions, whereas 6\% consistently allocate everything to the in-group member. Moving from complete in-group favoritism to equal treatment is associated with a 0.09 point increase in support for reducing domestic inequality and a 0.17 point increase in support for prioritizing the global over the domestic poor, both measured on four-point scales. These estimates are correlations rather than causal effects \cite[Table 1]{cappelen_universalism_2025}.

%===================================================================================================
%- allocation tasks, prioritization: Fabre et al 25, Fabre 26
%===================================================================================================

\citet{fabre_majority_2025} use a donation experiment in the United States and Europe to test support for international redistribution policies. Each respondent has to split a potential \$100 lottery prize between themselves and people living in poverty either in their home country or in Africa, with the recipient's location randomly assigned. In Europe, respondents allocate similar amounts to poor recipients in their own country and in Africa. In the United States, however, respondents give approximately \$3 more to poor recipients in their own country than to those in Africa. The estimated difference is \$2.509 and is statistically significant, whereas the corresponding difference in Europe is not statistically significant \citep[Figure S13 \& Table ED2]{fabre_majority_2025_si}.

In another allocation task, respondents have to allocate 100 points among six randomly selected policies. An equal allocation would therefore give approximately 16.7 points to each policy. On average, respondents in the U.S. and in Europe respectively allocated 21 and 20 points to a ``Global tax on millionaires'', placing it above this benchmark. By contrast, the policy ``Doubling foreign aid'' gathered 11 points, on average, in Europe and 9 in the U.S., placing it below the benchmark \citep[Figure S30]{fabre_majority_2025_si}.

The authors then ask respondents how many policies they support from a randomly assigned list. In the most comprehensive list, respondents evaluate four policies---the Global Climate Scheme, the National Redistribution Scheme, a coal exit, and ``marriage only for opposite-sex couples''---and can therefore report supporting between zero and four policies. The midpoint of this scale is two policies. The mean number of supported policies exceeds this midpoint in all four European countries: respondents support an average of 2.8 policies in France, 2.7 in Spain, 2.6 in the United Kingdom, and 2.2 in Germany. In the United States, the average is exactly at the midpoint, at 2.0 policies \citep[Figure S7]{fabre_majority_2025_si}.

In another allocation task, respondents have to allocate 100 points to different among six randomly selected policies. On average, respondents in the U.S. and in Europe respectively allocated 20 and 21 to a ``Global tax on millionaires''. While the policy ``Doubling foreign aid'' gathered 11 points, on average, in Europe and 9 in the U.S. \citep[Figure S30]{fabre_majority_2025_si}. 

\citet{fabre_majority_2025} asks respondents what share of the revenues from a worldwide tax on household wealth above \$5 million should be allocated to low-income countries. This share would be used to ``finance infrastructure and public services in low-income countries''. The median response is 30\% of the tax revenue across the five countries \citep[Figure S16 and ED4]{fabre_majority_2025_si}.

\citet{fabre_public_2026} conducts a revenue allocation task including five spending items---four of them concern domestic issues: the deficit, income tax, welfare programs, and education, while one concerns healthcare, education, and renewable energy in low-income countries. Respondents are asked how much of the revenue from a hypothetical wealth tax, applying a rate of 2\% to wealth above \$5 million, they would allocate to each item, after being informed of the revenue that the tax would generate in their country and in all low-income countries combined. Fabre finds, on average, that 17\% of the tax revenue would be allocated to global issues. The countries with the highest percentages are Spain and Saudi Arabia, both at 21\%, while the United States is at 17\% and Japan at 14\%.


%===================================================================================================
% - radical income redistribution: Fabre 26, unpublished results on top 5/8% in FR showing the limit of what could be accepted
%===================================================================================================

\citet{fabre_public_2026} examines support for ``a global tax on top incomes to finance poverty reduction in LICs''. Respondents are randomly presented with either a tax targeting the global top 1\% or a more progressive schedule covering the global top 3\%. Under the first proposal, after-tax income above \$120,000 per year at purchasing power parity is subject to an additional rate of 15\%. The second applies marginal rates of 15\%, 30\%, and 45\% above annual income thresholds of \$80,000, \$120,000, and \$1 million, respectively. The revenues would finance transfers associated with monthly poverty thresholds of \$250 under the first proposal and \$400 under the second, corresponding to approximately 2\% and 5\% of world nominal income. Absolute support for the tax on the global top 1\% reaches 56\% overall, 61\% in Europe, and 51\% in the United States \citep[Figure S40]{fabre_public_2026}. Among respondents who would personally pay either tax, 53\% accept the proposal overall; acceptance reaches 61\% in Europe, 32\% in Switzerland, and 79\% in Saudi Arabia. These affected respondents represent 6\% of the pooled sample, including 4\% in Europe, 11\% in Switzerland, and 21\% in Saudi Arabia \citep[Figure S41]{fabre_public_2026}.

\citet{fabre_public_2026} then asked respondents to construct their preferred global distribution of after-tax income by changing the shares of winners and losers and the redistribution from richer to poorer individuals. Respondents see both the current and proposed distributions: ``The current distribution is in red, and your custom one is in green'' and could preserve the status quo. The sliders' initial positions are randomly set to limit, and respondents are allowed to skip the task. The interactive exercise is available at \url{http://preferences-pol.fr/custom_global_redistr.html}, with an explanatory video at \url{https://youtu.be/gSfsQwczT7w}. Overall, 56\% of respondents declare themselves satisfied with their redistribution, whereas 43\% skip the task \citep[Figure S51]{fabre_public_2026}. Among respondents, the median proposal made 49\% of the world population better off and 18\% worse off, with a redistribution intensity of 5 out of 10 \citep[Figure S49]{fabre_public_2026}. This proposal would transfer 5.4\% of world income from richer to poorer individuals and raise the minimum monthly income to \$287 at purchasing power parity \citep[Figure S49]{fabre_public_2026}. Almost half of respondents (48\%) select a redistribution that would reduce their own income, compared with 9\% who would personally gain, while 10\% retain the current distribution. The average preferred redistribution would likewise transfer 5.4\% of world income, in this case from the top 27\% to the bottom 73\%, and establish a minimum monthly income of \$247 \citep[Figure S50]{fabre_public_2026}. The average and median proposals would reduce the global Gini coefficient from 0.67 to 0.59 and 0.58, respectively. At the individual level, the largest group of respondents (41\%) selected transfers representing between 4\% and 5\% of world income, while 55\% chose a redistribution producing a minimum monthly income between \$150 and \$350 \citep[Figure S56-57]{fabre_public_2026}. At the top of the distribution, the average proposal could be implemented through additional marginal tax rates of 7\% above \$25,000 and 16\% above \$40,000 per year at purchasing power parity \citep[Figure S69]{fabre_public_2026}.

Fabre finds original results from an online survey of 1,000 respondents representative of the French population conducted in 2026, showing that 41\% of respondents support a tax on the top 5\%, while 36\% oppose it, support decreases to 37\% for a tax on the top 8\%, while opposition increases to 39\%.

